%-----------------------------------------------------------------------------------------------------------------------------
% 1. Analýza problematiky                                           max. 15
% a. Prehľad súčasného stavu problematiky a teoretický rozbor       max. 10
\renewcommand{\bodyAnalyza}{} 
% - Do 10b – Študent preukázal široký prehľad v danej problematike, vie identifikovať a porovnať podobné riešenia, vie analyzovať dôvody a dopady riešenia, vie zhromaždiť a analyzovať potrebný teoreticky základ.   
% - Do 7b – Študent preukázal základný prehľad v danej problematike a analýza súčasného stavu je povrchná. 
% - Do 1b – Študent nepreukázal v dostatočnej miere, že má prehľad v súčasnom stave problematiky. 


% b. Práca s literárnymi zdrojmi                                    max. 5
\renewcommand{\bodyZdroje}{} 
% - Do 5b – Boli využité rôzne zdroje informácií (knižné, elektronické, vedecké články...). Pri rozbore alebo zmienke určitého faktu boli použité viaceré zdroje, ktoré podporujú dané tvrdenie. 
% - Do 3b – V práci bolo využité primerané množstvo rôznorodých zdrojov informácií, ale nie všetky podstatné tvrdenia sa opierali o konkrétny zdroj. 
% - Do 1b – Vzhľadom na problematiku je množstvo a zloženie zdrojov informácií nedostatočné. 

%------------------------------------------------------------------------------------------------------------------------------
% 2. Návrh riešenia - Fromulácia obmedzení, postupov a cieľov       max. 10
\renewcommand{\bodyNavrh}{} 
% - Do 10b – Návrh riešenia obsahuje všetky nasledovné zložky: stanovenie cieľov, opis možných prístupov vedúcich k dosiahnutiu cieľov; zamýšľaný postup riešenia; definície obmedzení, s ktorými riešenie uvažuje. 
% - Do 7b – Niektorá z uvedených zložiek chýba alebo sú niektoré časti návrhu nejasné či príliš stručné. 
% - Do 3b – Viaceré z uvedených zložiek chýbajú alebo sú nedostatočné, a teda z úvodnej časti práce nie je zrejmé, ako študent bude riešiť úlohy zadania a čo chce dosiahnuť. 

%------------------------------------------------------------------------------------------------------------------------------
% 3. Realizácia úloh                                                max. 20
% a. Názornosť pri dokumentovaní priebehu realizácie                max. 10
\renewcommand{\bodyRealizacia}{}  
% - Do 10b – Realizácia je dokumentovaná názorne. V práci sú použité náčrty, diagramy, blokové schémy, prípadne fotografie... 
% - Do 7b – Určité časti realizácie by si zaslúžili lepšiu názornosť. 
% - Do 3b – Opis realizácie si od čitateľa vyžaduje vysokú mieru predstavivosti. Nie je jasné ako boli realizované viaceré kroky. 

% b. Zdôvodnenie postupov a realizovaných krokov                    max. 5
\renewcommand{\bodyZdovodnenie}{}  
% - Do 5b – Všetky kroky a postupy pri riešení úloh zadania sú vhodne zdôvodnené a majú zjavné opodstatnenie. 
% - Do 3b – Niektoré kroky a postupy nie sú zdôvodnené a ani z kontextu nie je zrejmé ich opodstatnenie. 
% - Do 1b – Pri viacerých krokoch a postupoch nie je zrejmý dôvod, prečo boli vykonané. 

% c. Praktické využitie teoretických poznatkov                      max. 5                                 
\renewcommand{\bodyTeoria}{}    
% - Do 5b – Pri praktickej realizácii boli v dostatočnej miere využité teoretické poznatky. 
% - Do 3b – Kvalita dosiahnutých výsledkov je znížená z dôvodu nevyužitia niektorých potrebných teoretických poznatkov. 
% - Do 1b – Realizácia úloh nie je postavená na exaktných teoretických poznatkoch napriek tomu, že si to charakter úloh vyžaduje. 

%------------------------------------------------------------------------------------------------------------------------------
% 4. Prezentácia výsledkov                                          max. 20
% a. Funkčnosť riešenia a metodika jej overenia                     max. 5
\renewcommand{\bodyPrezentacia}{} 
% - Do 5b – Z práce je zrejmé, že výsledné riešenie je funkčné alebo spĺňa stanovené požiadavky. Spôsob overenia funkčnosti alebo plnenia požiadaviek je náležite opísaný.
% - Do 3b – Z kontextu nie je pochybnosť, že výsledné riešenie aspoň čiastočne spĺňa stanovené požiadavky, ale nie je uvedený spôsob ich overovania. 
% - Do 1b – Nie je poskytnutý žiaden dôkaz o tom, že výsledné riešenie je funkčné alebo spĺňa stanovené požiadavky. 

% b. Experimentálne overenie vlastností riešenia                    max. 10
\renewcommand{\bodyExperiment}{}
% - Do 10b – V práci sú zaznamenané a vhodne vizualizované údaje z dostatočného množstva experimentov (tabuľky, inžiniersky exaktné grafy...) s cieľom určiť vlastnosti a obmedzenia realizovaného riešenia. Jednotlivé experimenty sú náležite opísané. 
% - Do 7b – Bolo vykonané dostatočné množstvo experimentov, ale tie nie sú vhodne prezentované. 
% - Do 3b – Bolo vykonané nedostatočné množstvo experimentov pre určenie vlastností riešenia alebo kvalita experimentov neumožňuje vyvodiť relevantné závery. 

% c. Vlastné kvalitatívne zhodnotenie výsledkov                     max. 5
\renewcommand{\bodyZhodnotenie}{}  
% - Do 5b – Dosiahnuté výsledky a vlastnosti riešenia sú komplexne zhodnotené a okomentované v širších súvislostiach. 
% - Do 3b – Študent vyjadril svoj postoj k výsledkom a vykonal základnú analýzu vlastností riešenia. 
% - Do 1b – Zhodnotenie výsledkov je povrchné alebo nedostatočné. 

%------------------------------------------------------------------------------------------------------------------------------
% 5. Technická a programová dokumentácia (v prílohe, na proloženom médiu)   max. 10
\renewcommand{\bodyDokumentacia}{} 
% - Do 10b – Technická alebo programová dokumentácia dopĺňa prácu o všetky informácie potrebné na nezávislé zopakovanie postupov či experimentov a umožňuje používať a nadviazať na realizované riešenie. Je písaná prehľadne, štruktúrovane a názorne. 
% - Do 7b – Určité časti technickej dokumentácie nie sú spracované v dostatočnej kvalite na to, aby bolo možné postupy či experimenty presne zopakovať alebo dokumentácia nie je dostatočne prehľadná. 
% - Do 3b – Technická dokumentácia neposkytuje vhodné doplňujúce informácie. Na prácu nie je možné nadviazať a ani nie je možné zopakovať realizované postupy. 

%------------------------------------------------------------------------------------------------------------------------------
% 6. Celkové hodnotenie                                             max. 15
% a. Formálna úroveň                                                max. 5
\renewcommand{\bodyFormalna}{}
% - Do 5b – Bola vhodne zvolená štruktúra, rozsah aj formátovanie práce (napr. dodržaním odporúčanej šablóny). Rozloženie a členenie textu, veľkosť obrázkov, tabuliek a medzier medzi jednotlivými prvkami dokumentu pôsobí adekvátne. 
% - Do 3b – Určité prvky dokumentu sú z formálneho hľadiska použité neadekvátnym, rušivým spôsobom.
% - Do 1b – Práca po formálnej stránke nepôsobí profesionálne alebo seriózne.

% b. Úroveň písomného prejavu                                       max. 10
\renewcommand{\bodyPrejav}{} 
% - Do 10b – Odborné vyjadrovanie je presné a myšlienky sú podávané zrozumiteľne. Text nie je znehodnotený žiadnymi rušivými elementami (preklepy, gramatické či štylistické chyby). 
% - Do 7b - Odborné vyjadrovanie je v poriadku a myšlienky sú podávané zrozumiteľne, ale text je poznačený viacerými rušivými elementami (preklepy, gramatické či štylistické chyby). 
% - Do 5b – Odborné vyjadrovanie alebo zrozumiteľné podanie myšlienok má určité nedostatky, ale samotný text je bez rušivých elementov. 
% - Do 3b – V práci dominuje neodborné vyjadrovanie a niektoré vety či tvrdenia nedávajú zmysel. Častokrát nie je zrejmé, čo chcel autor daným tvrdením vyjadriť. 

% c. Náročnosť, pôvodnosť a využiteľnosť                            max. 10
\renewcommand{\bodyNarocnost}{}   
% - Do 10b – Spôsob riešenia úloh spĺňa požiadavku očakávanej náročnosti vyplývajúcej zo zadania. Práca vykazuje prekryv do 9% s inými autormi z CRZP. Výsledky sú aspoň čiastočne využiteľné aj v budúcich projektoch. 
% - Do 7b – Realizovaná práca nevykazuje takú náročnosť, aká sa predpokladá zo zadania alebo sú výsledky práce úplne nepoužiteľné v iných projektoch. Práca je však originálna. 
% - Do 4b – Práca vykazuje neoprávnený prekryv nad 25% s inými autormi z CRZP alebo náročnosť realizácie práce je veľmi nízka. 
